\section{Preliminaries}\label{sec:preliminaries}

\subsection{SPARQL Language}

\ac{RDF}~\cite{w3ConceptsAbstract} is a data model for representing knowledge on the Web.
Information is expressed as \emph{triples} of the form $(s, p, o)$, where $s$ is the \emph{subject} (the entity being described), $p$ is the \emph{predicate} (the property or relationship), and $o$ is the \emph{object} (the value or related entity).
A set of such triples forms an \ac{RDF} \ac{KG}, which can be thought of as a labeled directed graph: nodes are subjects and objects, and edges are predicates~\cite{w3ConceptsAbstract}.

\acp{IRI} (denoted $\mathcal{I}$) are globally unique identifiers used to name resources and properties.
Literals (denoted $\mathcal{L}$) are concrete data values such as strings, numbers, or dates.
Blank nodes (denoted $\mathcal{B}$) are anonymous resources with no global identifier.
A triple is formally defined as $(s, p, o)$ where $s \in \mathcal{I} \cup \mathcal{B}$, $p \in \mathcal{I}$, and $o \in \mathcal{I} \cup \mathcal{B} \cup \mathcal{L}$.
An \ac{RDF} \ac{KG} is a set of such triples.

SPARQL~\cite{w3SPARQLQuery} is the W3C standard language to query \ac{RDF} \ac{KG}.
SPARQL queries are matched against an \ac{RDF} graph by substituting variables for components of triples.
A \emph{triple pattern} extends a triple by allowing variables (drawn from a set $\mathcal{V}$, disjoint from $\mathcal{I}$, $\mathcal{B}$, and $\mathcal{L}$) in any position:
$\mathit{tp} = (s, p, o)$ where $s \in \mathcal{I} \cup \mathcal{B} \cup \mathcal{V}$, $p \in \mathcal{I} \cup \mathcal{V}$, and $o \in \mathcal{I} \cup \mathcal{B} \cup \mathcal{L} \cup \mathcal{V}$.

A \ac{BGP} is a set of triple patterns evaluated conjunctively, binding variables across all patterns.
A query built from \acp{BGP} is called a \ac{CQ}.
\acp{UCQ} are formed by taking the union of multiple \acp{GP} via a \ac{UGP}, allowing a query to match any of several alternative patterns.

Matching a graph pattern against an \ac{RDF} graph produces \emph{solution mappings}, partial functions from variables to \ac{RDF} terms.
The following definitions make precise two notions on solution mappings used throughout the paper.

\begin{definition}[Multiplicity]\label{def_multiplicity}
  Given a multiset or sequence $\Omega$ of solution mappings and a solution mapping $\mu$, we write $\operatorname{multiplicity}(\mu \mid \Omega)$ to denote the number of times $\mu$ appears in $\Omega$.
\end{definition}

\begin{definition}[Merge~\cite{w3ConceptsAbstract}]\label{def_merge}
  Two solution mappings $\mu_1$ and $\mu_2$ are \emph{compatible} if $\mu_1(v)$ and $\mu_2(v)$ are the same \ac{RDF} term~\cite{w3ConceptsAbstract} for every variable $v \in \operatorname{dom}(\mu_1) \cap \operatorname{dom}(\mu_2)$.
  The \emph{merge} of two compatible solution mappings is their union, written $\operatorname{merge}(\mu_1, \mu_2) = \mu_1 \cup \mu_2$.
\end{definition}

\subsection{Federated Queries}

A SPARQL federated query involves executing different segments of a query across multiple endpoints.
We denote the evaluation of a query over a \ac{KG} $G$ as $\eval{Q}{G}$, and the evaluation over a federation of endpoints $F$ as $\eval{Q}{F}$, where $F$ is defined as a set of federation members, formalized as a mapping $\mathit{fm}\colon \mathit{IRI} \mapsto G$.
Evaluating a query over $F$ involves decomposing $Q$ into subqueries $Q_i$, where each subquery $Q_i$ is evaluated on a knowledge graph $G_j \in \operatorname{codomain}(F)$ member of the federation.
We adopt a database-agnostic view of federation members: since query containment does not depend on the content of any particular knowledge graph, we formalize the notion of federation in terms of endpoint \acp{IRI} rather than data.

We distinguish three types of federated queries: \ac{UDF}, \ac{ESA}~\cite{Cheng2021}, and \ac{ASSF}.
In \ac{UDF}, the query specifies \texttt{SERVICE} clauses~\cite{w3SPARQLFederated} that explicitly name the target endpoint for each subquery; a variable may be used as the endpoint \ac{IRI}, in which case the query must satisfy \emph{Service-Safeness}~\cite{BuilAranda2013} to guarantee a finite federation.
In \ac{ESA}, a finite federation $F$ is provided before execution and each triple pattern is evaluated over every $G \in \operatorname{codomain}(F)$, with results accumulated via bag union~\cite{w3SPARQLFederated}.
In \ac{ASSF}, a finite $F$ is similarly given, but the query engine selects which federation member receives each triple pattern rather than broadcasting to all.

\begin{definition}[FedQPL operators~\cite{Cheng2021}]\label{def_fedqpl}
  We express federated query plans with three operators of the FedQPL formalism.
  The \emph{request} $\freq_{\mathit{fm}}^{P}$ evaluates a graph pattern $P$ against a single federation member $\mathit{fm}$.
  The \emph{multiway union} $\fmu$ and \emph{multiway join} $\fmj$ combine the results of their arguments by bag union and by join, respectively.
\end{definition}
